% 对象名字与伪代码关键字；变量和数学表达仍由原稿决定。
\floatname{algorithm}{算法}
\algrenewcommand{\algorithmicrequire}{\textbf{输入：}}
\algrenewcommand{\algorithmicensure}{\textbf{输出：}}
\algrenewcommand{\algorithmicfor}{\textbf{对于}}
\algrenewcommand{\algorithmicdo}{\textbf{执行}}
\algrenewcommand{\algorithmicif}{\textbf{若}}
\algrenewcommand{\algorithmicthen}{\textbf{则}}
\algrenewcommand{\algorithmicelse}{\textbf{否则}}
\algrenewcommand{\algorithmicend}{\textbf{结束}}
\algrenewcommand{\algorithmicreturn}{\textbf{返回}}
% 中文短标签以间距引出正文，取消 amsthm 默认的英文句点。
\newtheoremstyle{zhplain}{\topsep}{\topsep}{\itshape}{} {\bfseries}{}{0.5em}{}
\newtheoremstyle{zhdefinition}{\topsep}{\topsep}{}{} {\bfseries}{}{0.5em}{}
\theoremstyle{zhplain}
\newtheorem{theorem}{定理}
\newtheorem{lemma}[theorem]{引理}
\newtheorem{proposition}[theorem]{命题}
\theoremstyle{zhdefinition}
\newtheorem{definition}{定义}
\renewcommand{\proofname}{证明}
\makeatletter
% proof 带可选参数；实际正文保存在其内部命令中。
\expandafter\patchcmd\csname\string\proof\endcsname{\@addpunct{.}}{}{}{\PackageError{zh-objects}{Cannot remove proof heading punctuation}{Check the amsthm proof definition.}}
\makeatother
